\section{Introduction}\label{sec:introduction}

Let $S$ be a smooth integral complex projective surface and $L$ an ample
line bundle. We use additive notation for tensor powers and put $H=L^2$.
For a tuple of distinct points $\mathbf p=(p_1,\ldots,p_r)$, the multipoint
Seshadri constant is
\[
 \varepsilon(S,L;\mathbf p)
 =\inf_C\frac{L\cdot C}{\sum_{i=1}^r\mult_{p_i}C},
\]
where $C$ ranges over integral curves meeting at least one of the points.
If $\pi:Y\to S$ is the blow-up of these points, with exceptional curves
$E_1,\ldots,E_r$, the same constant is the largest $\lambda\ge0$ for
which $\pi^*L-\lambda\sum_iE_i$ is nef. Here a real divisor class is
\emph{nef} if it has nonnegative intersection with every integral curve.
This constant measures how much of the positivity of $L$ remains after
subtracting equal exceptional weights at the chosen points. Its universal
bound is
\begin{equation}\label{eq:upper-bound}
 \varepsilon(S,L;\mathbf p)\le\sqrt{\frac Hr};
\end{equation}
an elementary jet-count proof appears in Lemma~\ref{lem:sesha-upper}.

Write
\[
 U_r=\{(p_1,\ldots,p_r)\in S^r:p_i\ne p_j\text{ for }i\ne j\}.
\]
A property holds at \emph{very general} tuples if its failures are contained
in a countable union of proper Zariski-closed subsets of $U_r$.

\begin{theorem}\label{thm:main}
For every smooth integral complex projective surface $S$ and every ample
line bundle $L$ on $S$, there is an integer $r_0=r_0(S,L)$ with the following
property. For each integer $r\ge r_0$, there is a countable union
$Z_r$ of proper Zariski-closed subsets of $U_r$, with nonempty complement,
such that for every $\mathbf p\in U_r\setminus Z_r$ the class
\[
 \pi^*L-\sqrt{\frac{L^2}{r}}\sum_{i=1}^r E_i
\]
is nef on the blow-up at $\mathbf p$. Consequently,
\[
 \varepsilon(S,L;\mathbf p)=\sqrt{\frac{L^2}{r}}.
\]
\end{theorem}

Theorem~\ref{thm:main} resolves positively the qualitative Nagata--Biran
conjecture for surfaces, also called the qualitative
Nagata--Biran--Szemberg conjecture. The threshold is allowed to depend on
the polarized surface. The conclusion concerns the nef boundary class of
self-intersection zero. Thus it excludes every curve whose ratio of degree
to total point multiplicity falls below the numerical bound, even when
those multiplicities are unequal. The sharp threshold $r\ge10$ and strict nonhomogeneous
inequality in the plane Nagata problem at very general complex points
are a separate refinement, proved in
\cite[Theorem~1.1]{OpenAINagata2026}.
Corollary~\ref{cor:fields} extends eventual maximality to uncountable
algebraically closed fields of characteristic zero.

\subsection{Background}

Nagata formulated his plane conjecture in his work on Hilbert's
fourteenth problem and proved its strict multiplicity inequality when the
number of points is a square at least $16$
\cite{Nagata1959}; see also \cite[Chapter~III, Section~1, Proposition~1]{Nagata1965}. Multipoint Seshadri constants place that
problem in the geometry of arbitrary polarized surfaces. The qualitative
formulation asks for equality in \eqref{eq:upper-bound} for every
sufficiently large integer $r$; see
Syzdek--Szemberg~\cite[Conjecture~4.3]{SyzdekSzemberg2010}.
Strycharz-Szemberg--Szemberg~\cite[Conjecture~2.1]{StrycharzSzemberg2004}
formulate a stronger prediction specifying the threshold in terms of a
linear system on the surface.

Biran's stability theorem gives full symplectic packings for sufficiently
large numbers of equal balls, for rational symplectic classes on closed
four-manifolds \cite{Biran1999}; see also
\cite[Theorem~6.3]{Biran2001}. Buse--Hind--Opshtein subsequently proved
strong packing stability, allowing unequal sufficiently small balls, for
all closed symplectic four-manifolds
\cite[Theorem~1]{BuseHindOpshtein2016}. The algebraic statement asks for
nefness on blow-ups of the fixed complex surface. Its packing interpretation
uses K\"ahler forms, as made precise by
Eckl \cite[Theorem~0.5]{Eckl2017}; symplectic packing stability alone
does not fix the required complex structure.

On the algebraic side, Harbourne's estimates include maximality for all
sufficiently large $r$ with $rL^2$ a square, as well as asymptotic lower
bounds \cite[Theorem~I.1 and Proposition~I.2]{Harbourne2003}. The stated
results assume very ampleness; for an ample $L$, apply them to a very ample
power $aL$ and use $\varepsilon(S,aL;\mathbf p)=a\varepsilon(S,L;\mathbf p)$.
For these square values of $rL^2$, intersecting over $\C$ the countably
many open loci for bounds approaching the maximum gives equality at very
general tuples.
Assuming maximality for the corresponding plane constant, Ro\'e's transfer
gives the maximal bound when the surface
has a point with maximal one-point Seshadri constant
\cite[Theorem~3 and Corollary~7]{Roe2004}.
The multipoint product inequality, due to Biran for surfaces and stated
in greater generality by Ro\'e--Ross, propagates maximality from a fixed
configuration to suitable multiples of its size when the corresponding
plane constants are maximal
\cite[Theorem~1.1 and Remark~1.2]{RoeRoss2009}.
Our argument treats arbitrary ample polarizations and all integers beyond
one threshold directly.

There is also a distinction between eventual maximality and exact values
on particular surfaces. Dionne--Roth compute multipoint constants for
substantial families of polarizations on $\PP^1\times\PP^1$
\cite{DionneRoth2025}. The recent preprint of Laface--Ugaglia
\cite[Theorem~1]{LafaceUgaglia2026} claims the maximal irrational value
$1/\sqrt5$ for $\mathcal O(1,1)$ at ten very general points of that surface.
These fixed-surface and fixed-point-count results address a different
question from the uniform eventuality in Theorem~\ref{thm:main}.

\subsection{Proof overview}

The proof turns the problem into finite interpolation on the algebraic
torus $\T$. An \emph{injective jet test of order $m$} on a vector space
of functions means that only the zero function vanishes to order at least
$m$ at all the test points. We will show that, for every sufficiently
large $r$ and all positive integers $k,m$ with $m>k\sqrt{H/r}$, the
sections of $kL$ admit such a test at some $r$ distinct surface points.
The threshold in $r$ must work for arbitrarily large jet orders. This
uniformity is different from the asymptotic postulation theorem of
Alexander--Hirschowitz \cite[Theorem~1.1 and Corollary~1.2]{AlexanderHirschowitz2000},
where the degree threshold depends on a fixed bound for the point
multiplicities.

Section~\ref{sec:jets} establishes the transfer principle: a full-rank
jet matrix for the initial Taylor monomials remains full rank after a
sufficiently small analytic rescaling. We then choose one integral divisor
$V\in|dL|$ with an ordinary node.
In local coordinates its two branches are the coordinate axes. The
intersection number $kL\cdot V$ bounds the least Taylor exponent of a
section under weights $(1,t)$, after any copies of $V$ are removed.
The resulting exponents lie in a quadrilateral $kQ_t$.
A change of torus coordinates, followed by compression of the slices
parallel to $(1,1)$, gives a triangle $kP_t$ of area
$k^2H/2$ (Section~\ref{sec:node}). The slice compression has its own finite
jet-matrix limit. Applying these two transfers in succession carries an
injective torus test for this triangle to sections on the surface.
Monomial diagrams and nonzero interpolation minors are central in
Dumnicki's diagram method \cite[Proposition~13 and Theorem~14]{Dumnicki2007};
the particular weighted-exponent bound and diagonal compression needed
here are proved in full.

The next step uses the continuously varying shape of $P_t$. The equality
$\operatorname{area}(P_t)=H/2$ makes the torus jet threshold exactly
$\sqrt{H/r}$.
For every sufficiently large integer $r$, we tune this shape to a primitive
lattice direction and choose a subgroup of $\T$ with exactly $r$ points.
The subgroup is the kernel of two monomial equations with exponent
lattice of index $r$; it need not be a square grid. The parameter $t$ and
the subgroup are fixed before the jet orders $k,m$ are chosen.
Projection onto its characters reduces the jet test to a Laurent polynomial
supported in a triangle of area $(kw)^2/2$ and vertical range $kw$, where
$w=\sqrt{H/r}$. An elementary estimate using horizontal slices bounds its
multiplicity at the identity by $kw$ (Section~\ref{sec:lattice}).

Finally, the transferred surface tuple can depend on $k,m$, but injectivity
is a Zariski-open rank condition. Countably many such conditions give
simultaneous tests at very general tuples. To pass from equal vanishing
orders to all curves, suppose that a curve has negative intersection with
the square-zero boundary class. Subtracting a small multiple of that curve
and making a rational perturbation produces a positive-square divisor
beyond the boundary. Riemann--Roch makes a multiple effective; adding back
the curve yields a section with equal vanishing orders larger than the
critical slope, contradicting the jet tests (Section~\ref{sec:nef}).
